\documentclass[11pt]{article}

\usepackage[margin=1in]{geometry}
\usepackage{amsmath,amssymb,mathtools}
\usepackage{bm}
\usepackage{hyperref}
\usepackage{xcolor}
\usepackage{enumitem}
\usepackage{fancyvrb}

\hypersetup{
    colorlinks=true,
    linkcolor=blue,
    urlcolor=blue,
    citecolor=blue
}

\newcommand{\Tr}{\operatorname{Tr}}
\newcommand{\Id}{\mathbb{1}}
\newcommand{\Hphys}{\mathcal{H}_{\mathrm{phys}}}
\newcommand{\Haux}{\mathcal{H}_{\mathrm{aux}}}
\newcommand{\Htot}{\mathcal{H}_{\mathrm{phys}}\otimes\mathcal{H}_{\mathrm{aux}}}
\newcommand{\Sz}{S^z}
\newcommand{\Q}{Q}

\title{Purification, Thermal Ensembles, and \texttt{conserve="Sz"} in TenPy}
\author{}
\date{}

\begin{document}

\maketitle

\section*{Main conclusion}

For the J1-J2 triangular Heisenberg model in \texttt{finite\_T\_J1J2.py},
using
\[
\texttt{conserve="Sz"}
\]
does \emph{not} mean that the physical thermal ensemble is restricted to one
fixed magnetization sector. In the purification calculation used in the code,
the physical system is doubled by an auxiliary system. The U(1) charge constraint
is imposed on the \emph{doubled purified state}, where the auxiliary charges
compensate the physical charges. Therefore all physical total-\(\Sz\) sectors
remain present in the thermal trace.

The line
\[
\texttt{PurificationMPS.from\_infiniteT(...)}
\]
constructs the grand-canonical infinite-temperature purification. I checked the
installed TenPy source used in this environment, TenPy version \texttt{1.1.0}:
\texttt{PurificationMPS.from\_infiniteT} is the grand-canonical constructor,
whereas TenPy has a separate constructor for a fixed charge sector,
\[
\texttt{PurificationMPS.from\_infiniteT\_canonical(..., charge\_sector=...)}.
\]

Thus, for the Hamiltonian currently implemented,
\[
H
= \sum_{\langle i,j\rangle} J_1\, \bm S_i\cdot \bm S_j
+ \sum_{\langle\langle i,j\rangle\rangle} J_2\, \bm S_i\cdot \bm S_j
+ F_z \sum_i S_i^z,
\]
using \texttt{conserve="Sz"} is physically sensible and usually preferable.
It is a symmetry optimization. It is not a projection to the physical
\(S^z_{\mathrm{tot}}=0\), \(S^z_{\mathrm{tot}}=-N/2\), or any other fixed
magnetization sector.

\section{Thermal density matrices from the beginning}

Let the physical Hilbert space of \(N\) spin-\(1/2\) sites be
\[
\Hphys = \bigotimes_{i=1}^N \mathbb{C}^2.
\]
The thermal density matrix at inverse temperature \(\beta = 1/T\) is
\[
\rho_\beta
= \frac{e^{-\beta H}}{Z(\beta)},
\qquad
Z(\beta) = \Tr_{\Hphys}\left(e^{-\beta H}\right).
\]
For any physical observable \(O\),
\[
\langle O\rangle_\beta
= \Tr_{\Hphys}(\rho_\beta O)
= \frac{\Tr_{\Hphys}\left(e^{-\beta H}O\right)}
        {\Tr_{\Hphys}\left(e^{-\beta H}\right)}.
\]

If the Hamiltonian conserves total magnetization,
\[
\left[H, S^z_{\mathrm{tot}}\right]=0,
\qquad
S^z_{\mathrm{tot}} = \sum_i S_i^z,
\]
then the physical Hilbert space decomposes into charge sectors:
\[
\Hphys = \bigoplus_M \Hphys^{(M)},
\]
where \(M\) is the total \(S^z\). In this case the grand-canonical thermal
partition function, with no fixed magnetization constraint, is
\[
Z(\beta)
= \sum_M \Tr_{\Hphys^{(M)}}\left(e^{-\beta H_M}\right).
\]
Here \(H_M\) is the Hamiltonian restricted to sector \(M\).

A fixed-sector, or canonical-in-\(S^z\), calculation would instead use
\[
Z_M(\beta)
= \Tr_{\Hphys^{(M)}}\left(e^{-\beta H_M}\right)
\]
for a single chosen \(M\). This is \emph{not} what
\texttt{PurificationMPS.from\_infiniteT} does.

If \(F_z\neq 0\), then the Hamiltonian includes a longitudinal field
\[
F_z S^z_{\mathrm{tot}}.
\]
This biases the magnetization thermally, but it still does not fix
\(S^z_{\mathrm{tot}}\). The trace still sums over all \(M\), with field-dependent
Boltzmann weights.

\section{What purification does}

Purification rewrites a mixed density matrix on \(\Hphys\) as a pure state on a
larger Hilbert space
\[
\Htot = \Hphys \otimes \Haux.
\]
The auxiliary space \(\Haux\) is a copy of the physical Hilbert space.

Choose an orthonormal basis \(\{|n\rangle\}\) of \(\Hphys\). The infinite
temperature purification is the maximally entangled state
\[
|\Psi_{\infty}\rangle
= \frac{1}{\sqrt{\dim \Hphys}}
  \sum_n |n\rangle_{\mathrm{phys}}\otimes |n\rangle_{\mathrm{aux}}.
\]
Tracing out the auxiliary space gives
\[
\Tr_{\mathrm{aux}}\left(|\Psi_{\infty}\rangle\langle\Psi_{\infty}|\right)
= \frac{1}{\dim\Hphys}\Id_{\mathrm{phys}}.
\]
This is the infinite-temperature density matrix.

At finite temperature, imaginary-time evolution is applied only to the physical
system:
\[
|\Psi_\beta\rangle
= \left(e^{-\beta H/2}\otimes \Id_{\mathrm{aux}}\right)
|\Psi_{\infty}\rangle.
\]
Then
\[
\Tr_{\mathrm{aux}}\left(|\Psi_\beta\rangle\langle\Psi_\beta|\right)
\propto e^{-\beta H}.
\]
The normalization is
\[
\langle \Psi_\beta|\Psi_\beta\rangle
= \Tr_{\Hphys}\left(e^{-\beta H}\right)
= Z(\beta).
\]
Therefore
\[
\langle O\rangle_\beta
=
\frac{\langle\Psi_\beta|O|\Psi_\beta\rangle}
     {\langle\Psi_\beta|\Psi_\beta\rangle}.
\]

\section{One spin: the simplest example}

For one spin-\(1/2\), the physical basis is
\[
|\uparrow\rangle,\qquad |\downarrow\rangle.
\]
The infinite-temperature purification is
\[
|\Psi_\infty\rangle
= \frac{1}{\sqrt{2}}
\left(
|\uparrow\rangle_{\mathrm{phys}}|\uparrow\rangle_{\mathrm{aux}}
+
|\downarrow\rangle_{\mathrm{phys}}|\downarrow\rangle_{\mathrm{aux}}
\right).
\]
Tracing over the auxiliary spin gives
\[
\rho_\infty
= \Tr_{\mathrm{aux}}|\Psi_\infty\rangle\langle\Psi_\infty|
= \frac{1}{2}
\left(
|\uparrow\rangle\langle\uparrow|
+
|\downarrow\rangle\langle\downarrow|
\right)
= \frac{1}{2}\Id.
\]
Both physical \(S^z=+1/2\) and \(S^z=-1/2\) states are present. No physical
sector has been removed.

\section{Many spins: all physical magnetization sectors are included}

For \(N\) spins,
\[
|\Psi_\infty\rangle
= \frac{1}{\sqrt{2^N}}
\sum_{\sigma_1,\ldots,\sigma_N}
|\sigma_1,\ldots,\sigma_N\rangle_{\mathrm{phys}}
\otimes
|\sigma_1,\ldots,\sigma_N\rangle_{\mathrm{aux}}.
\]
The sum runs over all \(2^N\) spin configurations. These configurations include
every physical total magnetization sector:
\[
M=-\frac{N}{2},-\frac{N}{2}+1,\ldots,\frac{N}{2}.
\]
Thus the infinite-temperature state is not a fixed-\(M\) state of the physical
system. It is the full identity density matrix after tracing out the auxiliary
system.

\section{What \texttt{conserve="Sz"} means in an ordinary MPS}

In an ordinary pure-state MPS without purification, using a U(1) symmetry often
does mean that the represented physical wavefunction has a fixed total charge.
For spin systems, the charge is usually
\[
\Q = 2 S^z_{\mathrm{tot}}.
\]
For example, an ordinary MPS with fixed total charge \(\Q=0\) represents a pure
state in the physical \(S^z_{\mathrm{tot}}=0\) sector.

This intuition is the source of the confusion. Purification is different because
the MPS represents a vector in the doubled space
\[
\Hphys \otimes \Haux,
\]
not a wavefunction only in \(\Hphys\).

\section{What is the local MPS tensor?}

An MPS is just a factorized way to store the coefficients of a many-body state.
For an ordinary spin chain,
\[
|\psi\rangle
= \sum_{\sigma_1,\ldots,\sigma_N}
C_{\sigma_1,\ldots,\sigma_N}
|\sigma_1,\ldots,\sigma_N\rangle.
\]
The coefficient tensor \(C_{\sigma_1,\ldots,\sigma_N}\) is exponentially large.
An MPS rewrites it as a product of local tensors:
\[
C_{\sigma_1,\ldots,\sigma_N}
=
\sum_{\alpha_1,\ldots,\alpha_{N-1}}
B_1^{\sigma_1}{}_{\alpha_0,\alpha_1}
B_2^{\sigma_2}{}_{\alpha_1,\alpha_2}
\cdots
B_N^{\sigma_N}{}_{\alpha_{N-1},\alpha_N}.
\]
For finite boundary conditions, \(\alpha_0\) and \(\alpha_N\) have dimension one.
The indices \(\alpha_i\) are virtual or bond indices. The physical index
\(\sigma_i\) labels the local spin state.

In TenPy, the local MPS tensors are usually called \(B_i\). This \(B\) is not a
magnetic field; it is the conventional name for the local MPS tensor. Depending
on the canonical form, TenPy may internally store tensors in \(A\), \(B\), or
mixed form, but the conceptual object is the same: a local tensor with one
physical leg and two virtual legs.

For purification, each site has two local legs:
\[
p_i = \text{physical spin index}, \qquad
q_i = \text{auxiliary spin index}.
\]
Thus a purified MPS stores coefficients
\[
C_{p_1,q_1;\ldots;p_N,q_N}
\]
through local tensors
\[
B_i^{p_i,q_i}{}_{\alpha_{i-1},\alpha_i}.
\]
Equivalently,
\[
|\Psi\rangle
=
\sum_{\{p_i\},\{q_i\}}
C_{p_1,q_1;\ldots;p_N,q_N}
|p_1,\ldots,p_N\rangle_{\mathrm{phys}}
\otimes
|q_1,\ldots,q_N\rangle_{\mathrm{aux}}.
\]

\section{The infinite-temperature local tensor}

At infinite temperature, the purified state should be
\[
|\Psi_\infty\rangle
=
\frac{1}{\sqrt{2^N}}
\sum_{\sigma_1,\ldots,\sigma_N}
|\sigma_1,\ldots,\sigma_N\rangle_{\mathrm{phys}}
\otimes
|\sigma_1,\ldots,\sigma_N\rangle_{\mathrm{aux}}.
\]
This means that the physical and auxiliary configurations are identical in each
nonzero term:
\[
p_i=q_i \quad \text{for every site } i.
\]
Therefore the coefficient of a doubled configuration is
\[
C_{p_1,q_1;\ldots;p_N,q_N}
=
\frac{1}{\sqrt{2^N}}
\prod_{i=1}^N \delta_{p_i,q_i}.
\]

This is why the local infinite-temperature tensor is diagonal in \(p\) and \(q\):
\[
B_i^{p_i,q_i}
=
\frac{1}{\sqrt{2}}\delta_{p_i,q_i}
\]
up to the trivial left and right virtual legs present at \(\beta=0\).
For one site, explicitly,
\[
B^{\uparrow,\uparrow}=\frac{1}{\sqrt{2}},
\qquad
B^{\downarrow,\downarrow}=\frac{1}{\sqrt{2}},
\qquad
B^{\uparrow,\downarrow}=B^{\downarrow,\uparrow}=0.
\]
This is exactly the tensor version of
\[
|\Psi_\infty\rangle
=
\frac{1}{\sqrt{2}}
\left(
|\uparrow\rangle_{\mathrm{phys}}|\uparrow\rangle_{\mathrm{aux}}
+
|\downarrow\rangle_{\mathrm{phys}}|\downarrow\rangle_{\mathrm{aux}}
\right)
\]
on each site.

In the installed TenPy source, version \texttt{1.1.0}, the relevant file is

\begin{Verbatim}[fontsize=\small]
tenpy/networks/purification_mps.py
\end{Verbatim}

The relevant code is:

\begin{Verbatim}[fontsize=\small]
for i in range(L):
    p_leg = sites[i].leg
    B = npc.diag(1.0, p_leg, dtype, ['p', 'q']) / sites[i].dim ** 0.5
    # leg `q` has the physical leg with opposite `qconj`
    B = B.add_trivial_leg(0, label='vL', qconj=+1).add_trivial_leg(1, label='vR', qconj=-1)
    Bs[i] = B
\end{Verbatim}

There are two important facts in this block.

\begin{enumerate}[label=(\roman*)]
\item \texttt{npc.diag(..., ['p', 'q'])} makes a diagonal tensor between the
physical leg \(p\) and the auxiliary leg \(q\). This is the code version of
\[
B_i^{p_i,q_i}\propto\delta_{p_i,q_i}.
\]
\item The comment says that the auxiliary leg \(q\) has the physical leg with
opposite \texttt{qconj}. This is TenPy's way of saying that the auxiliary leg is
the dual, or bra-like, copy of the physical leg for charge bookkeeping.
\end{enumerate}

For the spin-\(1/2\) site itself, the relevant part of TenPy's
\texttt{SpinHalfSite} source is:

\begin{Verbatim}[fontsize=\small]
if conserve == 'Sz':
    chinfo = npc.ChargeInfo([1], ['2*Sz'])
    leg = npc.LegCharge.from_qflat(chinfo, [1, -1])
\end{Verbatim}

Thus the local physical charges are \(+1\) for \(|\uparrow\rangle\) and
\(-1\) for \(|\downarrow\rangle\), using the convention \(\Q=2S^z\).
The diagonal construction gives \(\delta_{p,q}\). The opposite \texttt{qconj}
on the auxiliary leg gives the charge cancellation explained next.

\section{What \texttt{conserve="Sz"} means in purification}

In TenPy's purification MPS, each site tensor has a physical leg and an auxiliary
leg. Call their local basis labels \(p\) and \(q\). For \texttt{SpinHalfSite}
with \texttt{conserve="Sz"}, the physical charges are
\[
\Q(|\uparrow\rangle)=+1,\qquad
\Q(|\downarrow\rangle)=-1,
\]
where \(\Q=2S^z\).

The important implementation detail is that the auxiliary leg \(q\) is used with
the opposite charge orientation. A nonzero tensor element obeys the charge
selection rule
\[
\Q(\alpha_i)
=
\Q(\alpha_{i-1}) + \Q(p_i)-\Q(q_i),
\]
up to TenPy's internal sign convention for incoming and outgoing legs. The
physical meaning is simple: the virtual bond charge tracks the accumulated
difference
\[
\Q_{\mathrm{bond}}(i)
=
\sum_{j\le i}\left[\Q(p_j)-\Q(q_j)\right].
\]
For a finite MPS with trivial boundary charges, the total doubled charge is
\[
\sum_{j=1}^N \left[\Q(p_j)-\Q(q_j)\right]=0.
\]

This is the charge that is conserved by the purification MPS with
\texttt{conserve="Sz"}:
\[
\Q_{\mathrm{doubled}}
=
\Q_{\mathrm{phys}}-\Q_{\mathrm{aux}}.
\]
It is \emph{not} the physical charge \(\Q_{\mathrm{phys}}\) alone.

At \(\beta=0\), the diagonal tensor \(B_{p,q}\propto\delta_{p,q}\) makes every
site locally neutral in this doubled sense:
\[
\Q(p_i)-\Q(q_i)=0.
\]
Thus all virtual bond charges are zero at infinite temperature. This does not
mean all physical spins are in the same \(S^z\) sector. It means that every
physical configuration is paired with an auxiliary configuration of the same
charge.

Because the auxiliary leg has the opposite effective charge, the physical charge
and auxiliary charge cancel:
\[
\Q_{\mathrm{phys}}(p) + \Q_{\mathrm{aux, effective}}(q) = 0
\quad\text{when }p=q.
\]
For example,
\[
|\uparrow\rangle_{\mathrm{phys}}|\uparrow\rangle_{\mathrm{aux}}
\quad\text{has}\quad
\Q_{\mathrm{phys}}=+1,\quad
\Q_{\mathrm{aux,effective}}=-1,
\]
so the doubled charge is zero. Similarly,
\[
|\downarrow\rangle_{\mathrm{phys}}|\downarrow\rangle_{\mathrm{aux}}
\quad\text{has}\quad
\Q_{\mathrm{phys}}=-1,\quad
\Q_{\mathrm{aux,effective}}=+1,
\]
so the doubled charge is also zero.

Therefore the purified MPS can have total doubled charge zero while still
containing all physical magnetization sectors.

Another way to say the same thing is:
\[
\texttt{conserve="Sz"}
\quad\Longrightarrow\quad
\text{TenPy conserves charge on the doubled purified MPS.}
\]
It does not imply
\[
\text{physical } S^z_{\mathrm{tot}} \text{ is fixed to one chosen value.}
\]
Only if one explicitly uses a canonical fixed-charge purification is a single
physical charge sector selected, for example
\[
\texttt{from\_infiniteT\_canonical}.
\]

\section{Why the auxiliary leg has opposite charge orientation}

The reason is that purification is closely related to vectorizing an operator.
Write a physical operator \(X\) as
\[
X=\sum_{p,q} X_{p,q}\, |p\rangle\langle q|.
\]
The index \(p\) is a ket index. The index \(q\) is a bra index. Under a U(1)
rotation generated by charge \(\Q\),
\[
|p\rangle \mapsto e^{i\theta \Q(p)}|p\rangle,
\qquad
\langle q| \mapsto e^{-i\theta \Q(q)}\langle q|.
\]
Therefore the operator basis element transforms as
\[
|p\rangle\langle q|
\mapsto
e^{i\theta[\Q(p)-\Q(q)]}
|p\rangle\langle q|.
\]
Its charge is not \(\Q(p)+\Q(q)\). Its charge is
\[
\Q(p)-\Q(q).
\]

Purification represents this operator-like object as a vector in a doubled
Hilbert space:
\[
|X\rangle\rangle
=
\sum_{p,q} X_{p,q}
|p\rangle_{\mathrm{phys}}
\otimes
|q\rangle_{\mathrm{aux}}.
\]
To preserve the operator charge \(\Q(p)-\Q(q)\), the auxiliary leg must carry
the opposite charge orientation. This is why the auxiliary leg is not treated as
a second independent physical spin with the same charge sign. It represents the
bra index of the density matrix.

The infinite-temperature operator is the identity,
\[
X=\Id=\sum_p |p\rangle\langle p|.
\]
Vectorizing it gives
\[
|\Id\rangle\rangle
=
\sum_p |p\rangle_{\mathrm{phys}}\otimes |p\rangle_{\mathrm{aux}}.
\]
Every term has doubled charge
\[
\Q(p)-\Q(p)=0.
\]
This is exactly why the local tensor is diagonal:
\[
B_{p,q}\propto\delta_{p,q}.
\]

\section{Why one still sets \texttt{conserve="Sz"} by hand}

The line in TenPy's purification code

\begin{Verbatim}[fontsize=\small]
# leg `q` has the physical leg with opposite `qconj`
\end{Verbatim}

does not by itself create a nontrivial \(S^z\) conservation law. It only says how
to orient the charge of the auxiliary leg relative to the physical leg. There
must already be a nontrivial charge on the physical leg.

That physical-leg charge is created by the site definition. In your model this
happens here:

\begin{Verbatim}[fontsize=\small]
site = SpinHalfSite(conserve=conserve)
\end{Verbatim}

If one passes

\begin{Verbatim}[fontsize=\small]
conserve="Sz"
\end{Verbatim}

then TenPy constructs the physical leg with charges
\[
\Q(|\uparrow\rangle)=+1,\qquad
\Q(|\downarrow\rangle)=-1.
\]
Then \texttt{from\_infiniteT} uses the same physical leg for \(p\), and the
opposite orientation for \(q\). The result is a nontrivial doubled U(1) symmetry:
\[
\Q_{\mathrm{doubled}}=\Q_{\mathrm{phys}}-\Q_{\mathrm{aux}}.
\]

If instead one passes

\begin{Verbatim}[fontsize=\small]
conserve=None
\end{Verbatim}

then \texttt{SpinHalfSite} uses a trivial charge leg. In that case, the
auxiliary leg still has the opposite \texttt{qconj}, but the charge itself is
trivial. Schematically,
\[
\Q(|\uparrow\rangle)=0,\qquad
\Q(|\downarrow\rangle)=0.
\]
Then
\[
\Q_{\mathrm{doubled}}=0-0=0
\]
for every tensor element, and the U(1) block structure is absent. The calculation
can still represent the same physics, but it stores tensors densely instead of
using the \(S^z\) block sparsity.

So there are two separate ingredients:

\begin{enumerate}[label=(\roman*)]
\item \texttt{conserve="Sz"} on \texttt{SpinHalfSite}: define a nontrivial U(1)
charge on the local physical basis.
\item opposite \texttt{qconj} on the purification auxiliary leg: make that
charge act as a bra/dual charge, so the purified tensor tracks
\(\Q_{\mathrm{phys}}-\Q_{\mathrm{aux}}\).
\end{enumerate}

Both are needed for a useful \(S^z\)-symmetric purification MPS. The first
supplies the charge labels. The second orients the auxiliary copy correctly.

This is why \texttt{conserve="Sz"} is not redundant. Without it, TenPy still
builds a grand-canonical infinite-temperature purification, but it does not use
the nontrivial \(S^z\) symmetry. With it, TenPy builds the same physical
grand-canonical purification in a symmetry-resolved tensor representation.

\section{Why TenPy's purification example does not use \texttt{conserve="Sz"}}

The official TenPy purification example uses the transverse-field Ising chain:

\begin{Verbatim}[fontsize=\small]
from tenpy.models.tf_ising import TFIChain

model_params = dict(L=L, J=1.0, g=1.2)
M = TFIChain(model_params)
psi = PurificationMPS.from_infiniteT(M.lat.mps_sites(), bc=bc)
\end{Verbatim}

This is a good minimal example of purification, but it is not an example of a
Hamiltonian with U(1) \(S^z\) conservation. The transverse field mixes
\(|\uparrow\rangle\) and \(|\downarrow\rangle\), so it changes total \(S^z\).
Equivalently,
\[
[H_{\mathrm{TFI}},S^z_{\mathrm{tot}}]\ne 0.
\]

Therefore one should not impose \texttt{conserve="Sz"} for that model. In fact,
TenPy's \texttt{SpinHalfSite} documentation says that with
\texttt{conserve="Sz"}, the onsite operators \texttt{Sx}, \texttt{Sy},
\texttt{Sigmax}, and \texttt{Sigmay} are excluded. That is exactly because they
carry nonzero \(S^z\) charge. A transverse-field Ising model needs such
spin-flip operators, so it must use a site representation without U(1)
\(S^z\) conservation, or at most a different symmetry such as a discrete parity
if the model supports it.

Thus the rule is not
\[
\text{``purification should avoid \texttt{conserve="Sz"}.''}
\]
The rule is
\[
\text{``use only symmetries that the Hamiltonian exactly preserves.''}
\]

For the J1-J2 Heisenberg Hamiltonian in this project,
\[
S_i^zS_j^z,\qquad S_i^+S_j^-,\qquad S_i^-S_j^+,\qquad F_zS_i^z
\]
all conserve total \(S^z\). Hence \texttt{conserve="Sz"} is appropriate for
that model.

\section{Why doubled charge conservation does not remove physical sectors}

The phrase ``the purified MPS has total doubled charge zero'' means
\[
\Q_{\mathrm{phys}}-\Q_{\mathrm{aux}}=0.
\]
It does not mean
\[
\Q_{\mathrm{phys}}=0.
\]
It means
\[
\Q_{\mathrm{phys}}=\Q_{\mathrm{aux}}.
\]
The common value can be anything.

For example, a term with all physical spins up and all auxiliary spins up has
\[
\Q_{\mathrm{phys}}=+N,\qquad \Q_{\mathrm{aux}}=+N,
\]
so
\[
\Q_{\mathrm{phys}}-\Q_{\mathrm{aux}}=0.
\]
This term is allowed. A term with all physical spins down and all auxiliary
spins down has
\[
\Q_{\mathrm{phys}}=-N,\qquad \Q_{\mathrm{aux}}=-N,
\]
so it is also allowed. Likewise every intermediate magnetization sector is
allowed.

Equivalently, the doubled-charge-zero space contains all operator blocks
\[
\Hphys^{(M)} \to \Hphys^{(M)}
\]
for every physical magnetization \(M\). It excludes only off-diagonal operator
blocks
\[
\Hphys^{(M)} \to \Hphys^{(M')}
\qquad (M\ne M').
\]
For a Hamiltonian conserving \(S^z_{\mathrm{tot}}\), the thermal operator
\[
e^{-\beta H}
\]
is block diagonal in \(M\). Hence the excluded off-diagonal blocks are exactly
zero anyway. Nothing physical is lost.

The partition function is still
\[
Z(\beta)
=
\sum_M \Tr_{\Hphys^{(M)}}e^{-\beta H_M}.
\]
All physical magnetization sectors contribute to the trace.

\section{Why physical observables are unaffected}

Let
\[
X_\beta=e^{-\beta H/2}.
\]
The purification is the vectorized operator
\[
|\Psi_\beta\rangle = |X_\beta\rangle\rangle.
\]
For a physical observable \(O\),
\[
\frac{\langle \Psi_\beta|O|\Psi_\beta\rangle}
     {\langle \Psi_\beta|\Psi_\beta\rangle}
=
\frac{\Tr(X_\beta^\dagger O X_\beta)}
     {\Tr(X_\beta^\dagger X_\beta)}
=
\frac{\Tr(e^{-\beta H}O)}
     {\Tr(e^{-\beta H})}.
\]
This is the ordinary thermal expectation value.

If \(H\) conserves \(S^z_{\mathrm{tot}}\), then \(X_\beta\) is block diagonal in
physical magnetization:
\[
X_\beta = \bigoplus_M e^{-\beta H_M/2}.
\]
The doubled-charge-zero purified MPS stores all of these blocks at once:
\[
|X_\beta\rangle\rangle
=
\bigoplus_M |e^{-\beta H_M/2}\rangle\rangle.
\]
The thermal trace still sums over \(M\). Thus the symmetry is not changing the
ensemble; it is only removing operator blocks that symmetry says are exactly
zero.

There is one useful consequence: if an observable changes total \(S^z\), such
as a single \(S^x_i\), its thermal expectation value is zero in a
U(1)-symmetric thermal state. TenPy also does not provide \texttt{Sx} or
\texttt{Sigmax} as onsite operators when \texttt{conserve="Sz"}, because those
operators carry nonzero charge. Charge-neutral observables such as \(H\),
\(S^z_i\), \(S_i^+S_j^-\), and combinations forming transverse correlations can
still be represented in symmetry-compatible ways.

\section{A useful table}

For \(N\) spins, consider a physical configuration with total physical charge
\[
\Q_{\mathrm{phys}} = 2S^z_{\mathrm{tot}} = M_Q.
\]
In the infinite-temperature purification, it is paired with the same spin
configuration on the auxiliary leg. Since the auxiliary leg has the opposite
effective charge, the doubled charge is
\[
\Q_{\mathrm{doubled}}
= \Q_{\mathrm{phys}} + \Q_{\mathrm{aux,effective}}
= M_Q - M_Q = 0.
\]

\[
\begin{array}{c|c|c|c}
\text{physical configuration} &
\Q_{\mathrm{phys}} &
\Q_{\mathrm{aux,effective}} &
\Q_{\mathrm{doubled}} \\
\hline
\text{all up} & +N & -N & 0 \\
\text{all down} & -N & +N & 0 \\
\text{half up, half down} & 0 & 0 & 0 \\
\text{any configuration} & M_Q & -M_Q & 0
\end{array}
\]

All physical magnetization sectors are compatible with the same total doubled
charge sector.

\section{Why imaginary-time evolution preserves the conclusion}

Your Hamiltonian is
\[
H
= \sum_{\langle i,j\rangle}
J_1
\left(
S_i^zS_j^z
+ \frac{1}{2}S_i^+S_j^-
+ \frac{1}{2}S_i^-S_j^+
\right)
+ \sum_{\langle\langle i,j\rangle\rangle}
J_2
\left(
S_i^zS_j^z
+ \frac{1}{2}S_i^+S_j^-
+ \frac{1}{2}S_i^-S_j^+
\right)
+ F_z\sum_i S_i^z.
\]
Each term conserves total \(S^z\):
\[
[H,S^z_{\mathrm{tot}}]=0.
\]
Therefore \(e^{-\beta H/2}\) also conserves physical \(S^z\). Starting from a
purification that contains all physical \(S^z\) sectors, imaginary-time
evolution evolves each sector independently:
\[
|\Psi_\beta\rangle
= \sum_M
\left(e^{-\beta H_M/2}\otimes \Id_{\mathrm{aux}}\right)
|\Psi_{\infty}^{(M)}\rangle.
\]
No sector is thrown away. The final thermal trace is still
\[
Z(\beta)=\sum_M \Tr_{\Hphys^{(M)}}e^{-\beta H_M}.
\]

\section{Mapping this to \texttt{finite\_T\_J1J2.py}}

The relevant code structure is:
\[
\texttt{site = SpinHalfSite(conserve=conserve)}
\]
inside \texttt{models/model\_J1J2.py}, and
\[
\texttt{psi = PurificationMPS.from\_infiniteT(M.lat.mps\_sites(), bc=bc)}
\]
inside \texttt{imag\_apply\_mpo}.

The first line says: use block-sparse tensors with a U(1) charge if
\texttt{conserve="Sz"}.

The second line says: initialize the grand-canonical infinite-temperature
purification.

The evolution operator is generated by
\[
\texttt{M.H\_MPO.make\_U(-d * dt, approx)}.
\]
Since \texttt{make\_U(t)} approximates \(e^{tH}\), using negative real
imaginary-time steps gives approximately
\[
e^{-dt H}.
\]
The purification state should evolve as
\[
|\Psi_\beta\rangle=e^{-\beta H/2}|\Psi_\infty\rangle.
\]
Thus if the state is evolved by \(dt\) in the ket, the physical density matrix
corresponds to an inverse-temperature increase
\[
\Delta\beta = 2dt.
\]
That is why the code has
\[
\texttt{beta += 2. * dt}.
\]

\section{What would be a fixed-sector calculation?}

A fixed-sector calculation would explicitly choose one physical total charge
sector. Examples:

\begin{enumerate}[label=(\alph*)]
\item Using TenPy's canonical purification constructor:
\[
\texttt{PurificationMPS.from\_infiniteT\_canonical(sites, charge\_sector=...)}
\]
This is designed for a canonical infinite-temperature ensemble in a specified
charge sector.

\item Using exact diagonalization with a charge-sector projection:
\[
\texttt{ExactDiag(model, charge\_sector=...)}
\]

\item Initializing a pure product state such as all spins down. That would not be
a thermal purification of the full Hilbert space.
\end{enumerate}

Your current code does none of these.

\section{When should \texttt{conserve="Sz"} not be used?}

Use \texttt{conserve="Sz"} only if the Hamiltonian conserves total \(S^z\).
The current J1-J2 Heisenberg model with longitudinal field does conserve it.

Do \emph{not} use \texttt{conserve="Sz"} if you add terms such as
\[
h_x\sum_i S_i^x,
\qquad
h_y\sum_i S_i^y,
\]
or any other term that changes total \(S^z\). In that case,
\[
[H,S^z_{\mathrm{tot}}]\neq 0,
\]
and the U(1) symmetry assumption is invalid. Then use
\[
\texttt{conserve=None}.
\]

\section{Why \texttt{conserve=None} and \texttt{conserve="Sz"} can differ numerically}

In exact arithmetic and with sufficiently large bond dimension, the two
calculations should give the same physical grand-canonical thermal observables
for this Hamiltonian:
\[
\texttt{conserve=None}
\quad\text{and}\quad
\texttt{conserve="Sz"}.
\]
However, with finite bond dimension and truncation, they can differ slightly.
Reasons include:

\begin{itemize}
\item The block-sparse calculation distributes the available bond dimension
across symmetry sectors.
\item The non-symmetric calculation has no sector structure, so truncation is
organized differently.
\item At the same nominal \(\chi_{\max}\), the effective variational spaces are
not identical.
\item Numerical roundoff and singular-value truncation can differ.
\end{itemize}

Such differences are truncation effects, not evidence that
\texttt{conserve="Sz"} computes a fixed physical magnetization ensemble.

\section{Practical checklist for the current code}

For the present J1-J2 triangular Heisenberg model:

\begin{itemize}
\item \texttt{conserve="Sz"} makes sense.
\item The ensemble initialized by \texttt{from\_infiniteT} is grand-canonical in
physical magnetization.
\item The physical trace includes all \(S^z_{\mathrm{tot}}\) sectors.
\item \texttt{conserve=None} is allowed but should be slower and usually less
efficient.
\item \texttt{conserve=None} is required only if the Hamiltonian breaks total
\(S^z\) conservation.
\end{itemize}

Two unrelated code-cleanup points:

\begin{itemize}
\item The loop \(\texttt{while beta < beta\_max}\) can overshoot by one step due
to floating-point roundoff. An integer step count is safer.
\item If the plotted \(C_v\) is meant to be specific heat per site, divide the
finite-difference result by \(N=L_xL_y\). Without that division, it is the total
heat capacity of the finite cluster.
\end{itemize}

\section*{Final statement}

The phrase \texttt{conserve="Sz"} is easy to misread. In this purification
context it should be read as:
\[
\text{``use \(S^z\) conservation as a tensor symmetry.''}
\]
It should not be read as:
\[
\text{``fix the physical thermal ensemble to one \(S^z\) sector.''}
\]

For the current J1-J2 Heisenberg Hamiltonian, using \texttt{conserve="Sz"} in
the purification calculation is physically correct.

\end{document}
